% -*- coding: utf-8 -*-
% !TEX program = xelatex
\documentclass[plain,noanswer,medmath]{../../template/jnuexam}
\usepackage{../../template/kysx}

\begin{document}


\examtitle{2014}{3}


\loadproblems{1}{2014P1}%载入试卷一
\loadproblems{2}{2014P2}%载入试卷二


\makepart{选择题}{1～8小题，每小题4分，共32分}


\begin{problem}
设$\lim_{n \to \infty} a_n = a$且$a \ne 0$，则当$n$充分大时有\pickout{}
\begin{abcd}
\item $|a_n| > \frac{|a|}{2}$．
\item $|a_n| < \frac{|a|}{2}$．
\item $a_n > a - \frac{1}{n}$．
\item $a_n > a + \frac{1}{n}$
\end{abcd}
\end{problem}


\useproblem{1}{1}{1}%【同试卷一第一(1)题】


\begin{problem}
设$p(x) = a + bx + cx^2 + \dx^3$，当$x \to 0$时，若$p(x) - \tan x$是比$x^3$高阶的无穷小，
则下列选项中错误的是\pickout{}
\begin{abcd}
\item $a = 0$．
\item $b = 1$，
\item $c = 0$．
\item $d = \frac{1}{6}$．
\end{abcd}
\end{problem}


\useproblem{1}{1}{2}%【同试卷一第一(2)题】


\useproblem{1}{1}{5}%【同试卷一第一(5)题】


\useproblem{1}{1}{6}%【同试卷一第一(6)题】


\useproblem{1}{1}{7}%【同试卷一第一(7)题】


\begin{problem}
设$X_1,X_2,X_3$为来自正态总体$N(0,\sigma^2)$的简单随机样本，
则统计量$\frac{X_1 - X_2}{\sqrt{2}|X_3|}$服从的分布为\pickout{}
\begin{abcd}
\item $F(1,1)$．
\item $F(2,1)$．
\item $t(1)$．
\item $t(2)$．
\end{abcd}
\end{problem}


\makepart{填空题}{9～14小题，每小题4分，共24分}


\begin{problem}
设某商品的需求函数为$Q = 40 - 2P$（$P$为商品价格），则该商品的边际收益为 \fillin{}．
\end{problem}


\begin{problem}
设$D$是由曲线$xy + 1 = 0$与直线$y + x = 0$及$y = 2$围成的有界区域，则$D$的面积为 \fillin{}．
\end{problem}


\begin{problem}
设$\int_0^a x\e^{2x} \dx = \frac{1}{4}$，则$a =$ \fillin{}．
\end{problem}


\begin{problem}
二次积分$\int_0^1 \dy \int_y^1 \left(\frac{\e^{x^2}}{x} - \e^{y^2} \right)\dx =$ \fillin{}．
\end{problem}


\useproblem{1}{2}{13}%【同试卷一第二(13)题】


\useproblem{1}{2}{14}%【同试卷一第二(14)题】


\makepart{解答题}{15～23小题，共94分}


\useproblem[points=10]{1}{3}{15}%【同试卷一第三(15)题】


\useproblem[points=10]{2}{3}{17}%【同试卷二第三(17)题】


\begin{problem}[points=10]%【类似试卷一第三(17)题】
设函数$f(u)$具有连续导数，$z = f(\e^x\cos y)$满足
$$\cos y\frac{\pdz}{\pdx} - \sin y\frac{\pdz}{\pdy} = (4z + \e^x\cos y) \e^{2x}.$$
若$f(0) = 0$，求$f(u)$的表达式．
\end{problem}


\begin{problem}[points=10]
求幂级数$\sum_{n = 0}^{\infty}(n + 1)(n + 3)x^n$的收敛域及和函数．
\end{problem}


\useproblem[points=10]{2}{3}{19}%【同试卷二第三(19)题】


\useproblem[points=11]{1}{3}{20}%【同试卷一第三(20)题】


\useproblem[points=11]{1}{3}{21}%【同试卷一第三(21)题】


\useproblem[points=11]{1}{3}{22}%【同试卷一第三(22)题】


\begin{problem}[points=11]
设随机变量$X$与$Y$的概率分布相同，$X$的概率分布为
$$P\{X = 0\} = \frac{1}{3},\quad P\{X = 1\} = \frac{2}{3},$$
且$X$与$Y$的相关系数$\rho_{XY} = \frac{1}{2}$．\par
\begin{enumerate*}
\item 求$(X,Y)$的概率分布；
\item 求$P\{X + Y \le 1\}$．
\end{enumerate*}
\end{problem}


\end{document}
